\ea\label{ex:text20-1}
KAVE GOROTO\\
\gll kave	goroto\\
     net	turtle\\
\glt `TURTLE NET'
\z

\ea\label{ex:text20-2}
O kave goroto vai Teapu na pakupaku raara bono naono to dao raara bono baneo.\\
\gll o	kave	goroto	vai	teapu	na	paku{$\sim$}paku	ra=ara	bono	naono	to	dao	ra=ara	bono	baneo\\
     \textsc{art}3.\textsc{sg}	net	turtle	\textsc{dem}6	Teop	\textsc{tm}2	\textsc{red}{$\sim$}make	1\textsc{pl}.\textsc{in}.\textsc{ipfv}=1\textsc{pl}.\textsc{in}.\textsc{pron}	\textsc{obj}.\textsc{art}3.\textsc{sg}	tree	\textsc{rel}	call	1\textsc{pl}.\textsc{in}.\textsc{ipfv}=1\textsc{pl}.\textsc{in}.\textsc{pron}	\textsc{obj}.\textsc{art}.3\textsc{sg}	breadfruit.tree\\
\glt `The Teop turtle net, we make it from a tree that we call the breadfruit tree.'
\z

\ea\label{ex:text20-3}
O baneo bona he toro sii vaamanae koa.\\
\gll o	baneo	bona	he	toro	sii	vaamanae	koa\\
     \textsc{art}3.\textsc{sg}	breadfruit.tree	\textsc{ana}.\textsc{sg}	\textsc{advers}	\textsc{tm}7	\textsc{dim}	quite.young	only\\
\glt `But this breadfruit tree must be quite young.'
\z

\ea\label{ex:text20-4}
A taba momohu toro paku ara, eara toro mosi gunaha o baneo.\\
\gll a	taba	momohu	toro	paku	ara	e-ara	toro	mosi	gunaha	o	baneo\\
     \textsc{art}2.\textsc{sg}	thing	first	\textsc{tm}7	do	1\textsc{pl}.\textsc{in}.\textsc{pron}	\textsc{h}-1\textsc{pl}.\textsc{in}.\textsc{pron}	\textsc{tm}7	cut	go.down	\textsc{art}3.\textsc{sg}	breadfruit.tree\\
\glt `The first thing we must do, we must cut down the breadfruit tree.'
\z

\ea\label{ex:text20-5}
Beve obete nana teo kasuana, eara repaa kahu a kapa nae bono kehaa to dao raara bono sui.\\
\gll beve	obete	nana	te=o	kasuana	e-ara	re=paa	kahu	a	kapa	nae	bono	kehaa	to	dao	ra=ara	bono	sui\\
     when.3\textsc{sg}.\textsc{pron}	lie	3\textsc{sg}.\textsc{ipfv}	\textsc{prep}=\textsc{art}3.\textsc{sg}	ground	\textsc{h}-1\textsc{pl}.\textsc{in}.\textsc{pron}	\textsc{consec}=\textsc{tm}3	scrape	\textsc{art}2.\textsc{sg}	bark	3\textsc{sg}.\textsc{poss}.\textsc{pron}	\textsc{obj}.\textsc{art}.3\textsc{sg}	shell	\textsc{rel}	call	1\textsc{pl}.\textsc{in}.\textsc{ipfv}=1\textsc{pl}.\textsc{in}.\textsc{pron}	\textsc{obj}.\textsc{art}.3\textsc{sg}	ark.clam\\
\glt `When it is lying on the ground, we scrape its bark off with a shell that we call \textit{sui} (`Ark Shellfish').'
\z

\ea\label{ex:text20-6}
Beara kahu vakavakavara e, eara repaa varoava e bono buaku ge o kukan o bon.\\
\gll be=ara	kahu	vaka-va-kavara	e	e-ara	re=paa	va-roava	e	bono	buaku	ge	o	kukan	o	bon\\
     when=1\textsc{pl}.\textsc{in}.\textsc{pron}	scrape	\textsc{red}{$\sim$}\textsc{advr}-complete	3\textsc{sg}.\textsc{pron}	\textsc{h}-1\textsc{pl}.\textsc{in}.\textsc{pron}	\textsc{consec}=\textsc{tm}3	\textsc{caus}-be.in.the.sun	3\textsc{sg}.\textsc{pron}	\textsc{obj}.\textsc{art}.3\textsc{sg}	two	or	\textsc{art}3.\textsc{sg}	three	\textsc{art}3.\textsc{sg}	day\\
\glt `When we have finished scraping it, we put it into the sun for two or three days.'
\z

\ea\label{ex:text20-7}
Beve raka vaantee, eara repaa hihi kahi e bono naono nae.\\
\gll beve	raka	va-antee	e-ara	re=paa	hihi	kahi	e	bono	naono	nae\\
     when.3\textsc{sg}.\textsc{pron}	dry	\textsc{advr}-enough	\textsc{h}-1\textsc{pl}.\textsc{in}.\textsc{pron}	\textsc{consec}=\textsc{tm}3	pull.off.fibres	from	3\textsc{sg}.\textsc{pron}	\textsc{obj}.\textsc{art}.3\textsc{sg}	wood	3\textsc{sg}.\textsc{pron}\\
\glt `When it is dry enough, we remove them  (the fibres) from its wood.'
\z

\ea\label{ex:text20-8}
Beara hihi vakavara e, eara repaa hupihupi e bono sui.\\
\gll be=ara	hihi	va-kavara	e	e-ara	re=paa	hupi{$\sim$}hupi	e	bono	sui\\
     when=1\textsc{pl}.\textsc{in}.\textsc{pron}	remove	\textsc{advr}-complete	3\textsc{sg}.\textsc{pron}	\textsc{h}-1\textsc{pl}.\textsc{in}.\textsc{pron}	\textsc{consec}=\textsc{tm}3	\textsc{red}{$\sim$}comb	3\textsc{sg}.\textsc{pron}	\textsc{obj}.\textsc{art}3.\textsc{sg}	ark.shell\\
\glt `When we have removed them, we comb them with the ark shell.'
\z

\ea\label{ex:text20-9}
Beara hupihupi vakavara e, a kiu vaa tea vinugu repaa taneo, ore tavus potee bona kanono.\\
\gll be=ara	hupi{$\sim$}hupi	va-kavara	e	a	kiu	vaa	tea	vinugu	re=paa	taneo	o=re	tavus	potee	bona	kanono\\
     when=1\textsc{pl}.\textsc{in}.\textsc{pron}	\textsc{red}{$\sim$}comb	\textsc{advr}-complete	3\textsc{sg}.\textsc{pron}	\textsc{art}2.\textsc{sg}	work	\textsc{lnk}	\textsc{compl}1	twist	\textsc{consec}=\textsc{tm}3	start	3.\textsc{pron}=\textsc{consec}	become	be.like	\textsc{obj}.\textsc{art}2.\textsc{sg}	string\\
\glt `When we have finished combing them (the fibres), the work of twisting starts, so that they become a string.'
\z

\newpage
\ea\label{ex:text20-10}
A kiu vai na antee nana bona buaku ge a dorana a otei tea peha kanono,\\
\gll a	kiu	vai	na	antee	nana	bona	buaku	ge	a	dorana	a	otei	te=a	peha	kanono\\
     \textsc{art}2.\textsc{sg}	work	\textsc{dem}6	\textsc{tm}2	need	3\textsc{sg}.\textsc{ipfv}	\textsc{obj}.\textsc{art}2.\textsc{sg}	two	or	\textsc{art}2.\textsc{sg}	four	\textsc{art}2.\textsc{sg}	man	\textsc{prep}=\textsc{art}2.\textsc{sg}	one2.\textsc{sg}	string\\
\glt `This work needs two or four men for one string,'
\z

\ea\label{ex:text20-11}
bea meha buaku ge a dorana vinugu batari bona meha kanono\\
\gll be=a	meha	buaku	ge	a	dorana	vinugu	bata=ri	bona	meha	kanono\\
     while=\textsc{art}2.\textsc{sg}	other2.\textsc{sg}	two	or	\textsc{art}2.\textsc{sg}	four	twist	\textsc{cont}=3\textsc{pl}.\textsc{ipfv}	\textsc{art}2.\textsc{sg}	other2.\textsc{sg}	string\\
\glt `while another two or four twist another string.'
\z

\ea\label{ex:text20-12}
A kiu vai toro vatotoena ni ara, tabae a kave toro bebeahu.\\
\gll a	kiu	vai	toro	va-totoena	ni	ara	tabae	a	kave	toro	bebeahu\\
     \textsc{art}2.\textsc{sg}	work	\textsc{dem}6	\textsc{tm}7	\textsc{mult}-share	\textsc{appl}	1\textsc{pl}.\textsc{in}.\textsc{pron}	because	\textsc{art}2.\textsc{sg}	net	\textsc{tm}7	long\\
\glt `This work we must share, because the net should become long.'
\z

\ea\label{ex:text20-13}
Bea kiu vaa tea vinugu kavara, eara repaa vavaantee nia maa kanono,\\
\gll be=a	kiu	vaa	tea	vinugu	kavara	e-ara	re=paa	va{$\sim$}va-antee	ni=a	maa	kanono\\
     when=\textsc{art}2.\textsc{sg}	work	\textsc{lnk}	\textsc{compl}1	twist	finished	\textsc{h}-1\textsc{pl}.\textsc{in}.\textsc{pron}	\textsc{consec}=\textsc{tm}3	\textsc{red}{$\sim$}\textsc{caus}-enough	\textsc{appl}=\textsc{art}2.\textsc{sg}	\textsc{plm}	string\\
\glt `When the work of twisting is finished, we have made the strings long enough,'
\z

\ea\label{ex:text20-14}
ei kou toro bebeahu vaantee.\\
\gll ei~kou	toro	bebeahu	va-antee\\
     because	\textsc{tm}7	long	\textsc{advr}-enough\\
\glt `because they must be long enough.'
\z

\ea\label{ex:text20-15}
Bea maa kanono vavaantee vakavara, eara repaa vaahana riori teo naono,\\
\gll be=a	maa	kanono	va{$\sim$}va-antee	va-kavara	e-ara	re=paa	vaa-hana	ri=ori	te=o	naono\\
     when=\textsc{art}2.\textsc{sg}	\textsc{plm}	string	\textsc{red}{$\sim$}\textsc{caus}-enough	\textsc{advr}-complete	\textsc{h}-1\textsc{pl}.\textsc{in}.\textsc{pron}	\textsc{consec}=\textsc{tm}3	\textsc{caus}-hang	3\textsc{pl}.\textsc{objm}=3\textsc{pl}.\textsc{pron}	\textsc{prep}=\textsc{art}3.\textsc{sg}	tree\\
\glt `When the strings have all become long enough, we hang them on a tree,'
\z

\ea\label{ex:text20-16}
eara repaa rahi vatotobin riori.\\
\gll e-ara	re=paa	rahi	va-totobin	ri=ori\\
     \textsc{h}-1\textsc{pl}.\textsc{in}.\textsc{pron}	\textsc{consec}=\textsc{tm}3	pull	\textsc{caus}-straight	3\textsc{pl}.\textsc{objm}=3\textsc{pl}.\textsc{pron}\\
\glt `then we pull them straight.'
\z

\ea\label{ex:text20-17}
Beve kavara o hum kiu bona, eara repaa taneo mana tea huha.\\
\gll beve	kavara	o	hum	kiu	bona	e-ara	re=paa	taneo	mana	tea	huha\\
     when.3\textsc{sg}.\textsc{pron}	finished	\textsc{art}3.\textsc{sg}	part	work	\textsc{ana}.\textsc{sg}	\textsc{h}-1\textsc{pl}.\textsc{in}.\textsc{pron}	\textsc{consec}=\textsc{tm}3	start	eventually	\textsc{compl}1	tie\\
\glt `When this part of the work is finished, we eventually start with the tying.'
\z

\ea\label{ex:text20-18}
A kave na tanetaneo raara bona taba to dao raara bona tavara.\\
\gll a	kave	na	tane{$\sim$}taneo	ra=ara	bona	taba	to	dao	ra=ara	bona	tavara\\
     \textsc{art}2.\textsc{sg}	net	\textsc{tm}2	\textsc{red}{$\sim$}start	1\textsc{pl}.\textsc{in}.\textsc{ipfv}=1\textsc{pl}.\textsc{in}.\textsc{pron}	\textsc{obj}.\textsc{art}2.\textsc{sg}	thing	\textsc{rel}	call	1\textsc{pl}.\textsc{in}.\textsc{ipfv}=1\textsc{pl}.\textsc{in}.\textsc{pron}	\textsc{obj}.\textsc{art}2.\textsc{sg}	tavara\\
\glt `The net, we start with a thing that we call \textit{tavara}.'
\z

\ea\label{ex:text20-19}
Evehee, beara mei rara tea taneo, eara toro varaarare ni roho a tatabonata na kave bara bebeahu nae.\\
\gll evehee	be=ara	mei	rara	tea	taneo	e-ara	toro	va-raa-rare	ni	roho	a	tatabonata	n=a	kave	bara	bebeahu	nae\\
     but	when=1\textsc{pl}.\textsc{in}.\textsc{pron}	do.not.yet	1\textsc{pl}.\textsc{in}.\textsc{ipfv}	\textsc{compl}1	start	\textsc{h}-1\textsc{pl}.\textsc{in}.\textsc{pron}	\textsc{tm}7	\textsc{mult}-\textsc{red}.sort.out	\textsc{appl}	before	\textsc{art}2.\textsc{sg}	width	3\textsc{sg}.\textsc{poss}=\textsc{art}2.\textsc{sg}	net	and2	length	3\textsc{sg}.\textsc{poss}.\textsc{pron}\\
\glt `But before we start, we must decide on the width of the net and its length.'
\z

\ea\label{ex:text20-20}
Beara vuha vaapeha meo ipinoo to paa vamoromoroko ni vuara,\\
\gll be=ara	vuha	vaa-peha	me=o	ipinoo	to	paa	va-moro{$\sim$}moroko	ni	vu=ara\\
     when=1\textsc{pl}.\textsc{in}.\textsc{pron}	breath	\textsc{caus}-one	for=\textsc{art}3.\textsc{sg}	measure	\textsc{rel}	\textsc{tm}3	\textsc{mult}-\textsc{red}{$\sim$}talk	\textsc{appl}	\textsc{pfv}=1\textsc{pl}.\textsc{in}.\textsc{pron}\\
\glt `When we have agreed on the size that we have talked about,'\footnote{lit. `when we are having one breath for the measure'}
\z

\ea\label{ex:text20-21}
eara paku vai a tavara to kahi ipinoo bata nom ara o mata kave.\\
\gll e-ara	paku	vai	a	tavara	to	kahi	ipinoo	bata	nom	ara	o	mata	kave\\
     \textsc{h}-1\textsc{pl}.\textsc{in}.\textsc{pron}	make	now	\textsc{art}2.\textsc{sg}	tavara	\textsc{rel}	\textsc{tm}1	measure	\textsc{cont}	\textsc{ipfv}	1\textsc{pl}.\textsc{in}.\textsc{pron}	\textsc{art}3.\textsc{sg}	mesh	net\\
\glt `we make the \textit{tavara} (a piece of bamboo) which we are going to use for measuring the meshes of the net.'
\z

\ea\label{ex:text20-22}
Be kavara amaa vamoromoroko, eara repaa taneo tea huha kave.\\
\gll be	kavara	a=maa	va-moro{$\sim$}moroko	e-ara	re=paa	taneo	tea	huha	kave\\
     when	finished	\textsc{art}2.\textsc{sg}=\textsc{plm}	\textsc{mult}-\textsc{red}{$\sim$}talk	\textsc{h}-1\textsc{pl}.\textsc{in}.\textsc{pron}	\textsc{consec}=\textsc{tm}3	start	\textsc{compl}1	tie	net\\
\glt `When the discussions are finished, we start tying the net.'
\z

\ea\label{ex:text20-23}
A kiu a beera, me kahi huha vuara bona maa bon amaa bero.\\
\gll a	kiu	a	beera	me	kahi	huha	vu=ara	bona	maa	bon	a=maa	bero\\
     \textsc{art}2.\textsc{sg}	work	\textsc{art}2.\textsc{sg}	big	and4	\textsc{tm}1	tie	\textsc{pfv}=1\textsc{pl}.\textsc{in}.\textsc{pron}	\textsc{obj}.\textsc{art}2.\textsc{sg}	\textsc{plm}	day	\textsc{art}2.\textsc{sg}=\textsc{plm}	many\\
\glt `The work is big, and we will tie (the net) many days.'
\z

\ea\label{ex:text20-24}
Be kavara a kiu vaa tea huha, eara repaa paku pete a kiu vaa tea ibun bono kehaa.\\
\gll be	kavara	a	kiu	vaa	tea	huha	e-ara	re=paa	paku	pete	a	kiu	vaa	tea	ibun	bono	kehaa\\
     when	finished	\textsc{art}2.\textsc{sg}	work	\textsc{lnk}	\textsc{compl}1	tie	\textsc{h}-1\textsc{pl}.\textsc{in}.\textsc{pron}	\textsc{consec}=\textsc{tm}3	do	subsequently	\textsc{art}2.\textsc{sg}	work	\textsc{lnk}	\textsc{compl}1	thread	\textsc{obj}.\textsc{art}.3\textsc{sg}	shell\\
\glt `When the work of tying is finsihed, we will continue with the work of threading the shells.'
\z

\ea\label{ex:text20-25}
Amaa kehaa bona to dao raara bono koriitoo.\\
\gll a=maa	kehaa	bona	to	dao	ra=ara	bono	koriitoo\\
     \textsc{art}2.\textsc{sg}=\textsc{plm}	shell	\textsc{ana}.\textsc{sg}	\textsc{rel}	call	1\textsc{pl}.\textsc{in}.\textsc{ipfv}=1\textsc{pl}.\textsc{in}.\textsc{pron}	\textsc{obj}.\textsc{art}.3\textsc{sg}	Tiger.Cowrie\\
\glt `These shells we call \textit{koriitoo} (`Tiger Cowrie').'
\z

\ea\label{ex:text20-26}
O koriitoo kou to kahi ibun raara teo roveana vaa paana kave, eove to kahi rahirahi vaadovo bata ni nom bona kave.\\
\gll o	koriitoo	kou	to	kahi	ibun	ra=ara	te=o	roveana	vaa	paa-n=a	kave	eove	to	kahi	rahi{$\sim$}rahi	va-dovo	bata	ni	nom	bona	kave\\
     \textsc{art}3.\textsc{sg}	Tiger.Cowrie	\textsc{emph}2	\textsc{rel}	\textsc{tm}1	thread	1\textsc{pl}.\textsc{in}.\textsc{ipfv}=1\textsc{pl}.\textsc{in}.\textsc{pron}	\textsc{prep}=\textsc{art}3.\textsc{sg}	seam	\textsc{lnk}	bottom-3\textsc{sg}.\textsc{poss}=\textsc{art}2.\textsc{sg}	net	\textsc{h}.3\textsc{sg}.\textsc{pron}	\textsc{rel}	\textsc{tm}1	\textsc{red}{$\sim$}pull	\textsc{caus}-sink	\textsc{cont}	\textsc{appl}	\textsc{ipfv}	\textsc{obj}.\textsc{art}2.\textsc{sg}	net\\
\glt `The Tiger Cowry shells we thread onto the seam at the bottom of the net, they will pull the net down (and) make it sink (into the deep sea).'
\z

\ea\label{ex:text20-27}
Be kavara a kiu bona, eara repaa gono pete maa o buo roravita mene vaataena kave.\\
\gll be	kavara	a	kiu	bona	e-ara	re=paa	gono	pete	maa	o	buo	roravita	mene	vaa-tae-n=a	kave\\
     when	finished	\textsc{art}2.\textsc{sg}	work	\textsc{ana}.\textsc{sg}	\textsc{h}-1\textsc{pl}.\textsc{in}.\textsc{pron}	\textsc{consec}=\textsc{tm}3	get	subsequently	hither	\textsc{art}3.\textsc{sg}	two3.\textsc{sg}	Helmet.Shell	for	\textsc{caus}-bottom-3\textsc{sg}.\textsc{poss}-\textsc{art}2.\textsc{sg}	net\\
\glt `When this work is finished, we get two Helmet Shells for the bottom of the net.'\footnote{lit. `for making the bottom of the net'}
\z

\ea\label{ex:text20-28}
O peho teo meho punete, o meho he teo meho punete.\\
\gll o	peho	te=o	meho	punete	o	meho	he	te=o	meho	punete\\
     \textsc{art}3.\textsc{sg}	one3.\textsc{sg}	\textsc{prep}=\textsc{art}3.\textsc{sg}	other3.\textsc{sg}	edge	\textsc{art}3.\textsc{sg}	other3.\textsc{sg}	\textsc{advers}	\textsc{prep}=\textsc{art}3.\textsc{sg}	other3.\textsc{sg}	edge\\
\glt `One for one edge, the other one for the other edge.'
\z

\newpage
\ea\label{ex:text20-29}
A bua kehaa bona to kahi ato vamasuu bata ni nom bona kave\\
\gll a	bua	kehaa	bona	to	kahi	ato	va-masuu	bata	ni	nom	bona	kave\\
     \textsc{art}2.\textsc{sg}	two2.\textsc{sg}	shell	\textsc{ana}.\textsc{sg}	\textsc{rel}	\textsc{tm}1	hold	\textsc{advr}-steady	\textsc{cont}	\textsc{appl}	\textsc{ipfv}	\textsc{obj}.\textsc{art}2.\textsc{sg}	net\\
\glt `These two shells will hold the net steady.'
\z

\ea\label{ex:text20-30}
to mene orooro kahi batana bona maa hum to kahi vavaagunaha kave raara.\\
\gll to	mene	oro{$\sim$}oro	kahi	bata-na	bona	maa	hum	to	kahi	va{$\sim$}vaa-gunaha	kave	ra=ara\\
     \textsc{rel}	\textsc{tm}6	\textsc{red}{$\sim$}float	from	\textsc{cont}-3\textsc{sg}.\textsc{ipfv}	\textsc{obj}.\textsc{art}2.\textsc{sg}	\textsc{plm}	place	\textsc{rel}	\textsc{tm}1	\textsc{mult}-\textsc{caus}-go.down	net	1\textsc{pl}.\textsc{in}.\textsc{ipfv}=1\textsc{pl}.\textsc{in}.\textsc{pron}\\
\glt `that otherwise would float away from the places where we would drop the net down.'
\z

\ea\label{ex:text20-31}
O buo kara haraharaa mee kahi pam ni raara teo buo roravita.\\
\gll o	buo	kara	hara{$\sim$}haraa	mee	kahi	pam	ni	ra=ara	te=o	buo	roravita\\
     \textsc{art}3.\textsc{sg}	two3.\textsc{sg}	line	\textsc{red}{$\sim$}feel	also	\textsc{tm}1	tie	\textsc{appl}	1\textsc{pl}.\textsc{in}.\textsc{ipfv}=1\textsc{pl}.\textsc{in}.\textsc{pron}	\textsc{prep}=\textsc{art}3.\textsc{sg}	two3.\textsc{sg}	Helmet.Shell\\
\glt `We tie two feeling lines to the Helmet Shells.'
\z

\ea\label{ex:text20-32}
O ogaoga na kave he na paku raara bono naono no baneo to paa hihi kahi vuara bona kapa nae.\\
\gll o	ogaoga	n=a	kave	he	na	paku	ra=ara	bono	naono	n=o	baneo	to	paa	hihi	kahi	vu=ara	bona	kapa	nae\\
     \textsc{art}2.\textsc{pl}	float	3\textsc{sg}.\textsc{poss}=\textsc{art}2.\textsc{sg}	net	\textsc{advers}	\textsc{tm}2	make	1\textsc{pl}.\textsc{in}.\textsc{ipfv}=1\textsc{pl}.\textsc{in}.\textsc{pron}	\textsc{obj}.\textsc{art}.3\textsc{sg}	wood	3\textsc{sg}.\textsc{poss}=\textsc{art}3.\textsc{sg}	breadfruit.tree	\textsc{rel}	\textsc{tm}3	remove	from	\textsc{pfv}=1\textsc{pl}.\textsc{in}.\textsc{pron}	\textsc{art}2.\textsc{sg}	bark	3\textsc{sg}.\textsc{poss}.\textsc{pron}\\
\glt `The floats of the net, however, we make from the wood of the breadfruit tree from which we have removed its bark.'
\z

\ea\label{ex:text20-33}
Na mosimosi vakakato raara beara ipinoo bata rae bona ta naono to paa vaantee vuara bona bebeahu no buo kobisi bebeahu na kuri raara.\\
\gll na	mosi{$\sim$}mosi	va-kakato	ra=ara	be=ara	ipinoo	bata	ra=e	bona	ta	naono	to	paa	va-antee	vu=ara	bona	bebeahu	n=o	buo	kobisi	bebeahu	n=a	kuri	raara\\
     \textsc{tm}2	\textsc{red}{$\sim$}cut	\textsc{caus}-short	1\textsc{pl}.\textsc{in}.\textsc{ipfv}=1\textsc{pl}.\textsc{in}.\textsc{pron}	when=1\textsc{pl}.\textsc{in}.\textsc{pron}	measure	\textsc{cont}	1\textsc{pl}.\textsc{in}.\textsc{ipfv}=3\textsc{sg}.\textsc{pron}	\textsc{art}2.\textsc{sg}	piece	wood	\textsc{rel}	\textsc{tm}3	\textsc{caus}-stop	\textsc{pfv}=1\textsc{pl}.\textsc{in}.\textsc{pron}	\textsc{art}2.\textsc{sg}	length	3\textsc{sg}.\textsc{poss}=\textsc{art}3.\textsc{sg}	two3.\textsc{sg}	finger	long	3\textsc{sg}.\textsc{poss}=\textsc{art}2.\textsc{sg}	hand	1\textsc{pl}.\textsc{in}.\textsc{poss}.\textsc{pron}\\
\glt `We cut it into short pieces, while measuring it with a piece of wood that we have made as long as the length of two middle fingers of our hand.'
\z

\ea\label{ex:text20-34}
Beara mosi kurukurus bata rae, a mosi toro kerun tamuana batana.\\
\gll be=ara	mosi	kuru{$\sim$}kurus	bata	ra=e	a	mosi	toro	kerun	tamuana	bata-na\\
     when=1\textsc{pl}.\textsc{in}.\textsc{pron}	cut	\textsc{red}{$\sim$}cut.off	\textsc{cont}	1\textsc{pl}.\textsc{in}.\textsc{ipfv}=3\textsc{sg}.\textsc{pron}	\textsc{art}2.\textsc{sg}	cut	\textsc{tm}7	diagonal	always	\textsc{cont}-3\textsc{sg}.\textsc{ipfv}\\
\glt `When we cut it into pieces, the cut must always be diagonal.'
\z

\ea\label{ex:text20-35}
Eara kahi mosi kurus bata nom, beo meho sii aba vaanana vaaraka bata ni rori bari teo suraa,\\
\gll e-ara	kahi	mosi	kurus	bata	nom	be=o	meho	sii	aba	vaa-nana	vaa-raka	bata	ni	rori	bari	te=o	suraa\\
     \textsc{h}-1\textsc{pl}.\textsc{in}.\textsc{pron}	\textsc{tm}1	cut	cut.off	\textsc{cont}	\textsc{ipfv}	while=\textsc{art}2.\textsc{pl}	other2.\textsc{pl}	\textsc{dim}	person	\textsc{caus}-warm	\textsc{caus}-dry	\textsc{cont}	\textsc{appl}	3\textsc{pl}.\textsc{ipfv}	4\textsc{pl}.\textsc{pron}	\textsc{prep}=\textsc{art}3.\textsc{sg}	fire\\
\glt `We will do the cutting, while other people are warming and drying the pieces of wood over the fire,'
\z

\ea\label{ex:text20-36}
Be kavara, eara repaa ibun mana riori tea kara kanono to dao raara bona roveana.\\
\gll be	kavara	e-ara	re=paa	ibun	mana	ri=ori	te=a	kara	kanono	to	dao	ra=ara	bona	roveana\\
     when	finished	\textsc{h}-1\textsc{pl}.\textsc{in}.\textsc{pron}	\textsc{consec}=\textsc{tm}3	thread	eventually	3\textsc{pl}.\textsc{objm}=3\textsc{pl}.\textsc{pron}	\textsc{prep}=\textsc{art}2.\textsc{sg}	string	string	\textsc{rel}	call	1\textsc{pl}.\textsc{in}.\textsc{ipfv}=1\textsc{pl}.\textsc{in}.\textsc{pron}	\textsc{obj}.\textsc{art}2.\textsc{sg}	seam\\
\glt `When (this is) finished, we thread them onto the string that we call the seam string.'
\z

\newpage
\ea\label{ex:text20-37}
O kanono bona he na pakupaku raara bona roge.\\
\gll o	kanono	bona	he	na	paku{$\sim$}paku	ra=ara	bona	roge\\
     \textsc{art}2.\textsc{pl}	string	\textsc{ana}.\textsc{sg}	\textsc{advers}	\textsc{tm}2	\textsc{red}{$\sim$}make	1\textsc{pl}.\textsc{in}.\textsc{ipfv}=1\textsc{pl}.\textsc{in}.\textsc{pron}	\textsc{obj}.\textsc{art}2.\textsc{sg}	coastal.hibiscus\\
\glt `This strings we make from the coastal hibiscus.'
\z
